\section{Foundation statements with synchronized scopes}\label{app:repairs}
The first two results retain the successful corrections from the English-v2 review. The chamber statement also supplies the smooth-readout qualification requested by the English-v3 referee. These statements concern the general foundation, not just the polynomial detector.

\subsection{Retained bias and constrained relative entropy}
\begin{proposition}[Biased constraints and positive-probability events]\label{prop:bias}
Let $P$ be a probability, $C$ a bounded vector observable and $B$ a bounded real observable. Put $Q^{\xi,B}=Z^{-1}e^{\xi\cdot C+B}P$, where $Z=P(e^{\xi\cdot C+B})$. If $Q^{\xi,B}C=c$, then for every $Q\ll P$ with $QC=c$,
\[
 D(Q\Vert P)-QB=D(Q\Vert Q^{\xi,B})+\xi\cdot c-\log Z.
\]
Thus $Q^{\xi,B}$ uniquely minimizes $D(Q\Vert P)-QB$ on this constraint class. It minimizes bare entropy when $B=0$, or when $QB$ is constant on the entire constraint class. If $P(A)>0$, the unique minimum-entropy law supported on $A$ is $P_A=P(\cdot\mid A)$, with
$D(Q\Vert P)=D(Q\Vert P_A)-\log P(A)$ there.
\end{proposition}
\begin{proof}
The tilted density is bounded above and below relative to $P$. Subtract its bounded logarithm $\xi\cdot C+B-\log Z$ from $\log(dQ/dP)$ and integrate. This proves the equality including infinite entropies, since only bounded quantities are subtracted. Nonnegativity of relative entropy, with equality exactly at equal probabilities, proves uniqueness and the special cases. On $A$, use $dP_A/dP=1/P(A)$ in the same calculation.
\end{proof}
For $P=(1/2,1/2)$, $C=0$ and $B=(0,\log3)$, the tilt $(1/4,3/4)$ is not the bare-entropy optimizer; it is the optimizer of the displayed biased objective. Conditioning on a zero-probability event requires a separately specified disintegration and is not covered by the event formula.

\subsection{Two density traces and explicit external flux}
Let a bounded domain be partitioned into finitely many moving smooth subdomains $D_\gamma(a)$. Internal interfaces have a normal from the minus branch to the plus branch and normal parameter velocity $v_\Sigma$. Assume $C^1$ local transports and branchwise $C^1$ extensions of $\rho_{a,\gamma},V_{a,\gamma},F_{a,\gamma}$, with bounds sufficient for differentiated volume integrals and boundary traces. On any moving external boundary, let $v_{\rm out}$ be velocity along the outward normal. Dots below are \emph{Eulerian parameter derivatives at fixed spatial position}, not material derivatives along the transport. Define
\[
 J_a(F)=\sum_\gamma\int_{D_\gamma(a)}\rho_{a,\gamma}e^{V_{a,\gamma}}F_{a,\gamma}\,dx.
\]
\begin{theorem}[Moving-domain response with both physical traces]\label{thm:traces}
Under these hypotheses,
\begin{align}
 J_a'(F)&=\sum_\gamma\int_{D_\gamma(a)}e^{V_\gamma}
 (\dot\rho_\gamma F_\gamma+\rho_\gamma\dot F_\gamma+
          \rho_\gamma F_\gamma\dot V_\gamma)\,dx\notag\\
 &\quad+\sum_\Sigma\int_\Sigma v_\Sigma
 (\rho_-e^{V_-}F_--\rho_+e^{V_+}F_+)\,dS\notag\\
 &\quad+\sum_{\Gamma\ {\rm external}}\int_\Gamma
              v_{\rm out}\rho_\gamma e^{V_\gamma}F_\gamma\,dS.
 \label{eq:two-traces}
\end{align}
A fixed external boundary has zero normal velocity. Equal internal density traces give the common-density special case. For $Z_a=J_a(1)>0$, the normalized derivative is
$\{J_a'(F)-(J_a(F)/Z_a)J_a'(1)\}/Z_a$.
\end{theorem}
\begin{proof}
Pull a branch integral to its fixed reference domain by $\chi_a$. For $f=\rho_\gamma e^{V_\gamma}F_\gamma$ and $w=\partial_a\chi_a\circ\chi_a^{-1}$, differentiating the Jacobian and integrand gives
$\int_{D_\gamma}(\partial_a f+\nabla f\cdot w+f\operatorname{div}w)$.
The last two terms are $\operatorname{div}(fw)$. The divergence theorem yields its boundary flux. Opposite outward normals on an internal interface combine the two fluxes as $v_\Sigma(f_--f_+)$. External faces contribute their stated outward flux, not an omitted term. Expanding $\partial_af$ and then applying the quotient rule proves the result.
\end{proof}
For example $J(a)=\int_0^a\psi^2+2\int_a^1\psi^2$ has derivative $-\psi(a)^2$, exactly the two-trace term. A smooth common ambient density is a legitimate special case, not a reason to replace two unequal physical traces by one.

\subsection{Trajectory, readout and preparation regularity are distinct}
A regular finite-event chamber has a fixed itinerary of at most $M$ events on common compact neighborhoods. Flights, guards and resets are $C^{r+1}$, competing guards and endpoint events have positive separation, and each encountered guard satisfies
$|g_t+D_zg\,f^-|\ge\kappa>0$. The finite horizon is fixed.

\begin{theorem}[Chamber response with a specified smooth reporting map]\label{thm:chamber-repair}
In such a chamber, event times, one-sided event states and terminal \emph{states} are jointly $C^r$ in physical parameter and initial state, with uniform bounds. A report has this regularity when it is specified by a jointly $C^r$ map of these quantities and the parameter, with bounded derivatives on the common compact set. Arbitrary Borel readouts are not asserted to be smooth.

For integrated response, take either a fixed finite preparation or explicitly
\[
 \mu_a=Z_a^{-1}(X_a)_\#(w_a\lambda),\qquad
 Z_a=\int w_a\,d\lambda\ge z_*>0,
\]
where $\lambda$ is fixed and finite, $w_a\ge0$, and $X_a$ takes values in the common chamber. Weights and transports are jointly measurable and $C^r$ in the parameter almost everywhere in their source coordinate. Require common integrable envelopes for all derivative products through order $r$ of the composed weighted tests. Then the preparation integrals are $C^r$, with transport, weight, observable and normalizer derivatives retained.

For distributional response fix the reporting space to be a $d$-dimensional torus, or a fixed compact smooth $d$-manifold equipped with its Sobolev scale, and use the specified smooth report map into that space. Analogous integrable envelopes for weighted report jets give a strongly $C^r$ report law in $H^{-s}$ for $s>d/2+r$. Finite tags use a direct sum of fixed reporting spaces. A fixed ambient embedding may instead be used, with the Sobolev condition stated in that ambient dimension.
\end{theorem}
\begin{proof}
Smooth flight dependence follows from the variational ordinary differential equation. The implicit function theorem at an event applies because the guard time derivative is bounded away from zero. If $s^-$ is incoming sensitivity at nominal event time, differentiation gives
\[
 \tau'=-\frac{g_a+D_zg\,s^-}{g_t+D_zg\,f^-},\qquad
 s^+=D_z\mathcal R\,s^-+\mathcal R_a+
 (D_z\mathcal R f^-+\mathcal R_t-f^+)\tau'.
\]
The last correction compares states at the same physical time. Repeated implicit differentiation solves the highest event-time derivative by division by the same nonzero guard derivative; every other term involves known lower jets. Compose with resets and flights at most $M$ times. Compact bounds and guard separation give joint $C^r$ states and uniform derivatives without an unnoticed itinerary change. Composition with the explicitly smooth reporting map proves the report assertion, not an assertion about arbitrary measurable maps.

Let $Y_a(z)$ denote this report starting from $X_a(z)$. For a test $F_a$ the unnormalized integral is $J_a=\int w_a(z)F_a(Y_a(z))\lambda(dz)$. Its first derivative is
\[
 J_a'=\int\{\dot w_aF_a(Y_a)+w_a\dot F_a(Y_a)
                       +w_a DF_a(Y_a)\dot Y_a\}\,d\lambda,
\]
where $\dot Y_a$ includes the initial-transport derivative. Finite Leibniz and chain rules give all higher derivatives. The assumed envelopes justify differentiation and continuity by dominated convergence. Differentiate $J_a/Z_a$ with the quotient recursion, retaining every derivative of $Z_a$.

For laws on the torus, the $j$th derivative of $y\mapsto\delta_y$ has squared $H^{-s}$ norm bounded by a constant times
$\sum_{k\in\mathbb Z^d}(1+|k|^2)^{-s}|k|^{2j}$, finite for $s>d/2+r$. Fourier dominated convergence gives strong $C^r$ dependence. A finite smooth chart/partition argument proves the corresponding statement on a fixed compact manifold. Apply the chain rule and the jet envelopes to the Bochner integral $\int w_a\delta_{Y_a}\,d\lambda$ and then normalize. The fixed-preparation case is parameter-independent transport and weights.
\end{proof}

The smooth no-event system $\dot x=a$, $x(0)=0$, $T=1$, with fixed Borel report $\one_{\{x(T)>0\}}$, produces $\one_{\{a>0\}}$. It violates the smooth-report hypothesis, not the trajectory conclusion. Similarly, the cusp-weight preparation $(1/2+|a|)\delta_0+(1/2-|a|)\delta_1$ violates the weight hypothesis, not smooth identity dynamics. The cartridge theorem instead integrates a fixed ambient preparation through its moving insertion guard using the complete collision-tube calculation. It does not apply a regular-chamber estimate uniformly at grazing.
